% SPDX-License-Identifier: CC-BY-4.0
% Copyright (c) 2026 Martin Schröder <info@swedishembedded.com>

\section{Related work}\label{sec:related}

Only works that were read are cited.

\paragraph{Reinforcement learning on a speech model.}
HappyRobot describes tuning a text-to-speech model with reinforcement learning: several takes of
a sentence are sampled, each is scored by a multimodal judge's naturalness and by word error rate,
and the model is moved toward the better takes with group-relative policy optimisation, a KL penalty
to the original voice and a human check on what is shipped \citep{happyrobot2026tts}. The post
reports no quantitative results and gives no account of how a user's own corrections enter. Group-relative
advantages come from \citet{shao2024deepseekmath}; preference optimisation without a reward model
from \citet{rafailov2023dpo}. We use these as baselines and as the optimiser, and add the lesson layer
that decides what a correction is and where it goes.

\paragraph{Pronunciation control.}
Markup standards give a synthesizer explicit pronunciations: SSML phoneme and lexicon elements
\citep{w3c2010ssml} and the Pronunciation Lexicon Specification \citep{w3c2008pls}. Our lexicon lessons
are the same idea learned from a conversation rather than authored, with provenance, an owner and a
rollback.

\paragraph{Teaching an agent.}
This paper makes no comparison with work on agent memory, learning from feedback or editing a model's knowledge: none was read for it, so none is cited.
